% =====================================================================
%  Hard versus fuzzy eye-tracking gating of EEG windows: a controlled simulation study
%  ALL EXPERIMENTAL RESULTS IN THIS MANUSCRIPT ARE OBTAINED ON SYNTHETIC DATA.
%  Tables/figures in generated/ are produced by gfer/analyze_sim.py and gfer/analyze_noise.py.
% =====================================================================
\documentclass[11pt,a4paper]{article}
\usepackage[T1]{fontenc}
\usepackage[utf8]{inputenc}
\usepackage{lmodern}
\usepackage[margin=2.3cm]{geometry}
\usepackage{amsmath,amssymb,bm}
\usepackage{graphicx}
\usepackage{booktabs,multirow,tabularx,array}
\usepackage{xcolor}
\usepackage{algorithm}
\usepackage{algpseudocode}
\usepackage{tikz}
\usetikzlibrary{positioning,arrows.meta,fit,backgrounds,calc}
\usepackage[hidelinks]{hyperref}
\usepackage{caption}
\usepackage{framed}
\captionsetup{font=small,labelfont=bf,labelsep=space}
\newcommand{\pending}[1]{\textcolor{red!75!black}{\textbf{[#1]}}}
\newcommand{\simnote}{\textsuperscript{$\ast$}}

\title{Hard versus Fuzzy Eye-Tracking Gating of EEG Windows in a Sugeno Fuzzy Rule-Mining Emotion Recognizer: A Controlled Simulation Study}
\author{\pending{Author names and affiliations}}
\date{}

\begin{document}
\maketitle

% =====================================================================
\begin{abstract}
\noindent
Multimodal emotion recognition from electroencephalography (EEG) usually assumes that a stimulus label holds throughout the stimulus, although emotional responses are intermittent. A recent method addresses this by using an eye-tracking feature as a hard gate that retains only EEG segments at which the feature is high; the retained segments are described by wavelet-band entropy features and classified with genetic-algorithm-tuned Sugeno fuzzy rules. This paper asks whether replacing the crisp gate with a fuzzy gate, in which Low, Medium and High memberships of the normalised eye-tracking feature are mapped by a zero-order Sugeno rule base to a continuous window weight, is a better way to use the eye-tracking evidence. The weights enter feature selection, rule mining, genetic-algorithm fitness and trial-level aggregation, while all other components are unchanged. Because the true affective relevance of an EEG window cannot be observed in real recordings, the comparison is carried out in a controlled simulation: two synthetic 62-channel EEG and eye-tracking datasets (5 subjects, 60 one-minute trials, four classes each) with known emotional episodes, evaluated with trial-disjoint and leave-one-subject-out cross-validation over three seeds and five wavelet bands. \textbf{All results are obtained on synthetic data and are not estimates of real-world EEG performance.} When the gate feature tracked the episodes (pupil diameter, AUC 0.99), gating raised mean trial accuracy from 35.8\% to 55.2\% (hard) and 54.9\% (fuzzy) on the harder dataset and from 49.1\% to 66.0\% and 65.7\% on the easier one ($p\leq0.005$ against no gate). With the pupil gate, the fuzzy and hard gates did not differ significantly under either protocol ($p>0.4$), the fuzzy gate did not outperform a simple linear soft gate, and under increasing gate-signal noise the fuzzy gate did not degrade more gracefully than the hard gate. When the gate feature was chosen by minimum-redundancy-maximum-relevance ranking, as in the original method, fixation features were selected that carried little episode information (AUC 0.57) and neither gate helped consistently. Within this simulation, the form of the gate matters far less than whether the gating signal identifies the relevant windows.
\end{abstract}

\noindent\textbf{Keywords} Electroencephalogram (EEG) $\cdot$ Eye tracking $\cdot$ Emotion recognition $\cdot$ Fuzzy gating $\cdot$ Sugeno fuzzy inference system $\cdot$ Simulation study

\begin{framed}
\noindent\small\textbf{Data notice.} Every experimental number, table and figure in Sect.~\ref{sec:results} reports results on \emph{synthetic} EEG and eye-tracking data generated by the procedure in Sect.~\ref{sec:data}. Results marked \simnote{} are simulation results. They characterise the behaviour of the gating mechanisms under known ground truth and must not be compared with accuracies reported on real EEG corpora.
\end{framed}

% =====================================================================
\section{Introduction}
\label{sec:intro}

Automatic recognition of human affective states supports applications in mental-health assessment, adaptive human--computer interaction and user-experience evaluation. Among the physiological signals used for this purpose, the electroencephalogram (EEG) is attractive because it records cortical activity non-invasively with millisecond temporal resolution and at a cost far below that of functional neuroimaging \cite{Alarcao2019,Jenke2014}. EEG-based emotion recognition has progressed from band-power features with conventional classifiers to entropy-based and deep representations, and accuracies above 80\% are routinely reported on benchmark corpora such as DEAP, MAHNOB-HCI and the SEED family \cite{Koelstra2012,Soleymani2012db,ZhengLu2015,Zheng2019}.

Eye tracking provides a complementary, low-cost view of the same affective process. Pupil diameter increases with emotional arousal irrespective of valence \cite{Bradley2008,Partala2003}, and fixation and saccade statistics reflect the allocation of attention to a stimulus \cite{Lim2020}. Fusing EEG with eye-movement features improves recognition over either modality alone, for example from 78.51\% (EEG) and 77.80\% (eye movements) to 87.59\% with fuzzy-integral decision fusion \cite{Lu2015}, and to 85.11\% for four emotions on SEED-IV with multimodal deep learning \cite{Zheng2019}.

A less frequently addressed problem is the \emph{temporal validity of labels}. Emotion corpora assign one label to an entire stimulus, yet there is no guarantee that the intended emotion is present throughout the stimulus or that it is elicited at a predictable moment. When a long recording is cut into short windows that all inherit the stimulus label, many windows carry little affective information. Farhadi Sedehi et al.\ \cite{FarhadiSedehi2024a} proposed to use eye tracking as a \emph{gate} for this purpose: eight eye-tracking features are ranked with minimum-redundancy-maximum-relevance (mRMR) selection, the top-ranked feature (mean vertical pupil diameter on their data) is compared with a threshold, and only EEG segments at the corresponding moments are analysed with discrete wavelet transform (DWT) band entropies, mRMR feature selection and a Sugeno fuzzy inference system (S-FIS) whose rules are refined by a genetic algorithm (GA). They reported 94.42\% and 85.20\% alpha-band accuracy on SEED-IV and MAHNOB-HCI \cite{FarhadiSedehi2024a}, and later reused pupil-based gating in a Granger-causality model \cite{FarhadiSedehi2024b}.

The hard-threshold form of the gate has properties that motivate investigation. A window marginally below the threshold is discarded while one marginally above is retained with full weight, so small measurement fluctuations change the evidence used for training and testing. Pupil diameter is dominated by the light reflex and modulated by cognitive load, gaze position and tracking quality \cite{Mathot2018,Raiturkar2016,Holmqvist2012}, so treating a threshold crossing as a certain indicator of affective relevance ignores substantial uncertainty. Fuzzy set theory \cite{Zadeh1965} offers a natural way to express graded relevance, soft gating has a long history in machine learning \cite{Jacobs1991,Arevalo2020}, and fuzzy memberships have been used to weight training samples \cite{LinWang2002}. To the best of our knowledge, however, fuzzy memberships of eye-tracking features have not been used to weight EEG windows for emotion recognition.

Evaluating such a gate on real data faces a basic obstacle: the ground-truth affective relevance of each EEG window is unobservable, so one cannot tell whether a gate retains the right windows or merely changes the classifier's input distribution. This paper therefore studies the question in a \emph{controlled simulation} in which emotional episodes are known by construction. The research question is: \emph{Within an otherwise unchanged S-FIS-GA pipeline, does a fuzzy eye-tracking gate that assigns each EEG window a continuous relevance weight provide better or more robust recognition than the hard-threshold gate?} This was posed as a hypothesis to be tested, not assumed.

\subsection{Motivation}
Graded relevance is, in principle, a more faithful representation of an uncertain ocular indicator than a crisp threshold; a fuzzy gate never leaves a trial without evidence; and, because the classifier is itself a fuzzy rule base, a fuzzy gate keeps the whole pipeline interpretable. Whether these properties translate into accuracy or robustness is an empirical question that a simulation with known ground truth can answer cleanly before costly real-data studies.

\subsection{Contributions}
\begin{enumerate}
  \item We formalise the hard eye-tracking gate of \cite{FarhadiSedehi2024a} as a crisp window-weighting function and propose a fuzzy gate (M1) based on a Low/Medium/High partition and a zero-order Sugeno rule base, of which the hard gate (M0) is a limiting case.
  \item We integrate the window weights consistently into weighted mRMR selection, weighted entropy discretisation, weighted look-up-table rule mining, weighted GA fitness and weighted trial-level aggregation, so that M0 and M1 differ only in the gate.
  \item We compare M0, M1, a no-gate reference and a non-fuzzy linear soft gate on two synthetic datasets with known episodes, under trial-disjoint and leave-one-subject-out cross-validation, with paired statistical tests, a gate-noise robustness sweep and a direct measurement of how well each gate locates the episodes.
  \item We find\simnote{} that fuzzy and hard gating perform equivalently within the precision of the experiment, that the fuzzy gate is not more robust to gate noise, and that the choice of gating \emph{signal} dominates the choice of gating \emph{function}; in particular, label-driven mRMR selection of the gate feature can select features that do not indicate affective relevance.
\end{enumerate}

\subsection{Organization of the paper}
Section~\ref{sec:related} reviews related work and Sect.~\ref{sec:gap} states the research gap. Section~\ref{sec:method} describes the two gating models and the shared pipeline. Section~\ref{sec:data} describes the simulation design and Sect.~\ref{sec:setup} the experimental protocol. Section~\ref{sec:results} presents and discusses the results, Sect.~\ref{sec:limits} the limitations, and Sect.~\ref{sec:conclusion} concludes. Table~\ref{tab:abbr} lists the abbreviations.

\begin{table}[t]
\centering\small
\caption{List of abbreviations}\label{tab:abbr}
\begin{tabular}{@{}ll@{\hspace{2em}}ll@{}}
\toprule
Abbreviation & Description & Abbreviation & Description\\
\midrule
EEG & Electroencephalogram & mRMR & Min.-redundancy max.-relevance\\
ET & Eye tracking & S-FIS & Sugeno fuzzy inference system\\
DWT & Discrete wavelet transform & GA & Genetic algorithm\\
SampEn & Sample entropy & MF & Membership function\\
PermEn & Permutation entropy & LOSO & Leave-one-subject-out\\
SymbEn & Symbolic dynamics entropy & CV & Cross-validation\\
PhEn & Phase entropy & AUC & Area under the ROC curve\\
M0 & Hard-gated baseline & M1 & Proposed fuzzy-gated model\\
A1 & No-gate reference & A4 & Linear soft-gate ablation\\
SIM-H / SIM-E & Synthetic datasets (harder / easier) & pp & Percentage points\\
\bottomrule
\end{tabular}
\end{table}

% =====================================================================
